% !TeX root = main.tex
% !TEX encoding = UTF-8 Unicode
% Ensure this file is saved with UTF-8 encoding

% Problem Statement, Background, and Motivation
% The need for research

\chapter{Introduction}
\label{Chapter1}

This thesis presents the design, implementation, and evaluation of a \acrshort{riscv} \acrfull{soc} featuring a specialized \acrfull{dnn} accelerator, conducted within the \acrfull{ceslab} of the University of Science, VNUHCM (\acrshort{hcmus}). This work represents a foundational step in our laboratory's strategic transition to the emerging \acrshort{riscv} ecosystem and open-source \acrshort{soc} design automation solutions for future chip development and fabrication.

\section{Background and Motivation}

The rapid advancement of artificial intelligence and machine learning has created an unprecedented demand for specialized computing architectures. \acrlongpl{dnn}, which form the backbone of modern \acrshort{ai} applications, require enormous computational resources that traditional general-purpose processors struggle to provide efficiently. This computational challenge, combined with the exponential growth in model complexity and the need for energy-efficient solutions, has driven the development of custom hardware accelerators specifically designed for \acrshort{dnn} workloads.

Traditionally, the \acrshort{ceslab} has been involved in designing and developing hardware accelerators (or \acrshort{ip} Cores) on \acrshort{fpga} using \acrshort{intel} \acrshort{quartus}, particularly leveraging the Platform Designer tool with Nios soft processors and Avalon Bus interfaces. While this approach has served well for \acrshort{fpga}-based prototyping, it presents significant limitations for our laboratory's strategic evolution toward chip fabrication and \acrshort{vlsi} design. The Avalon Bus, though functional, lacks the sophistication and industry adoption of modern interconnect standards such as \acrshort{axi4} and TileLink. More critically, the \acrshort{intel} \acrshort{quartus} ecosystem is primarily optimized for \acrshort{fpga} implementation and lacks the comprehensive \acrshort{vlsi} design flow capabilities essential for \acrshort{asic} development and chip fabrication.

As \acrshort{ceslab} envisions a future centered on chip prototyping and semiconductor fabrication, the limitations of our current toolchain have become increasingly apparent. The proprietary nature of \acrshort{intel}'s tools not only imposes licensing costs but also restricts our ability to customize the design flow, integrate with modern \acrshort{vlsi} tools, and contribute to the broader open-source hardware community. Furthermore, the transition from \acrshort{fpga} prototypes to actual silicon requires a design methodology that seamlessly bridges both domains—a capability that traditional \acrshort{fpga}-centric tools cannot provide.

The emergence of the \acrshort{riscv} instruction set architecture has fundamentally changed the landscape of processor design and \acrshort{soc} development, offering a pathway to address these limitations. As an open-source \acrshort{isa}, \acrshort{riscv} eliminates the barriers imposed by proprietary architectures, providing unprecedented flexibility for customization and modification. More importantly, this openness has catalyzed the development of comprehensive design ecosystems that support the entire \acrshort{soc} development flow, from high-level specification through \acrshort{fpga} prototyping to final \acrshort{vlsi} implementation and chip fabrication.

\section{A Survey of Modern Open-Source SoC Frameworks}

To address the limitations of our previous proprietary toolchain and select a platform aligned with \acrshort{ceslab}'s strategic goal of \acrshort{vlsi} development, we conducted a survey of prominent open-source \acrfull{soc} design frameworks. The landscape is rich with powerful, mature options, each with a distinct design philosophy and target application domain. Key alternatives considered include the LiteX, PULP, and OpenTitan platforms.

\subsection{LiteX: The Flexible SoC Builder}

LiteX is a versatile framework that utilizes Migen, a \acrshort{python}-based hardware description language, to facilitate the rapid construction of custom \acrshortpl{soc} \cite{litex_github,litex_paper}. Its primary strength lies in its flexibility, offering a vast library of general-purpose \acrshort{ip} and support for over 150 \acrshort{fpga} boards from a wide range of vendors \cite{litex_boards_repo}. While LiteX provides an exceptional environment for rapid \acrshort{fpga} prototyping and experimentation, its toolchain is not inherently optimized for a direct, end-to-end \acrshort{vlsi} design flow in the same way as other frameworks. For our laboratory's specific goal of a seamless transition from \acrshort{fpga} prototyping to \acrshort{asic} implementation, LiteX presented a less direct path.

\subsection{PULP Platform: The Low-Power Parallel Computing Specialist}

The PULP (Parallel Ultra-Low-Power) Platform, developed at \acrshort{ethzurich}, is a leader in energy-efficient, parallel computing for edge \acrshort{ai} and \acrshort{iot} applications \cite{pulp_platform_org}. The platform is built around a family of highly-optimized \acrshort{riscv} cores and has been validated through an impressive 51 silicon tape-outs, demonstrating its maturity and industrial relevance \cite{pulp_implementation}. However, the entire PULP architecture is meticulously optimized for its specific domain of ultra-low-power, multi-core data processing. Our project required a general-purpose \acrshort{soc} architecture capable of integrating a large, tightly-coupled, memory-intensive \acrshort{dnn} accelerator, a different architectural paradigm than PULP's primary focus.

\subsection{OpenTitan: The Security-First Root of Trust}

OpenTitan is a high-profile, open-source project managed by lowRISC to create a transparent and verifiable silicon "root of trust" (RoT) \cite{opentitan_org}. The framework, backed by a coalition including Google, provides a comprehensive suite of security-hardened \acrshort{ip}, such as cryptographic accelerators and secure boot managers, built around the Ibex \acrshort{riscv} core \cite{opentitan_github}. While OpenTitan represents a landmark achievement in open-source hardware security, its purpose is highly specialized. It is engineered to provide verifiable security, not high-performance computation. As our thesis is focused on \acrshort{dnn} acceleration research, the OpenTitan framework was not a suitable foundation.

\section{Chipyard: the framework with Full-Stack Integration Capability}

Among the various \acrshort{riscv} \acrshort{soc} design frameworks, Chipyard has emerged as the most comprehensive and mature solution for academic research and industrial \acrshort{vlsi} development. Developed at \acrshort{ucb}, Chipyard provides a complete, open-source \acrshort{soc} design environment that spans from high-level hardware generators written in \acrshort{chisel} to full \acrshort{vlsi} implementation flows optimized for \acrshort{asic} fabrication. This end-to-end capability is precisely what \acrshort{ceslab} requires to transition from \acrshort{fpga}-based prototyping to actual chip development.

The framework's advantages over traditional \acrshort{fpga} design flows is evident in several key areas. Unlike the proprietary Avalon Bus used in \acrshort{intel}'s ecosystem, Chipyard leverages industry-standard interconnects including TileLink and \acrshort{axi4}. \acrshort{axi4}, in particular, has become the de facto standard for modern \acrshort{soc} design due to its superior performance characteristics, extensive tool support, and widespread adoption across the semiconductor industry. This adherence to industry standards ensures that designs developed using Chipyard can seamlessly integrate with commercial \acrshort{ip} cores and industry-standard \acrshort{vlsi} tools.

Most critically, Chipyard includes a fully integrated, open-source \acrshort{vlsi} flow that encompasses synthesis, place-and-route, and timing closure using industry-standard tools such as Cadence Genus and Innovus. This \acrshort{vlsi}-optimized design flow enables the direct translation of high-level hardware descriptions into fabrication-ready layouts, supporting multiple process technologies including advanced nodes from foundries like \acrshort{tsmc} and \acrshort{globalfoundries}. This capability is essential for \acrshort{ceslab}'s strategic goal of developing chip prototypes and eventual commercial silicon products.

A key innovation of Chipyard is its emphasis on \textit{full-stack integration}. Unlike traditional approaches that design and evaluate accelerators in isolation, Chipyard enables the creation of complete, \acrshort{linux}-capable systems that capture real-world effects such as operating system overhead, memory hierarchy interactions, and system-level resource contention. This holistic approach is crucial for understanding the true performance characteristics of specialized accelerators when deployed in practical computing environments—insights that are invaluable when transitioning from \acrshort{fpga} prototypes to optimized \acrshort{asic} implementations.

The framework includes the Gemmini accelerator generator, an open-source \acrshort{dnn} accelerator specifically designed for systolic array-based matrix computations. Gemmini embodies the principles of efficient \acrshort{dnn} acceleration through its configurable systolic array architecture, which maximizes data reuse and minimizes the costly movement of data between computational units and external memory—a fundamental challenge known as the "memory wall" in computer architecture.

\section{Core Selection within Chipyard: A Comparison of Rocket and BOOM}

A key advantage of the Chipyard framework is that it is a generator capable of instantiating different processor cores. This presented a critical design choice for our system: selecting the \acrshort{cpu} core that best aligned with the project's objectives. The two primary candidates developed at \acrshort{ucb} and available within Chipyard are the Rocket Core and the \acrfull{boom}.

\begin{itemize}
    \item The Rocket Core is a mature, stable, and widely-used 5-stage, single-issue, in-order processor core. It is highly configurable and provides a solid foundation for general-purpose computing, capable of booting full operating systems like \acrshort{linux}. Its relative simplicity and extensive documentation make it a reliable and robust choice for systems where the primary focus is on integrating other components, such as custom accelerators \cite{asanovic2016rocketchip}.
    \item The \acrfull{boom} is a synthesizable, parameterized, out-of-order \acrshort{riscv} core designed for high-performance applications \cite{celio2016boom_spec, celio_phdthesis_2018boom}. It features a much more complex microarchitecture than Rocket, with a deeper pipeline and sophisticated hardware for dynamic scheduling, register renaming, and speculative execution. \acrshort{boom} is designed to compete with high-performance commercial cores, targeting performance levels suitable for desktop or server-class applications \cite{zhao2020sonicboom_ieee}.
\end{itemize}

For this thesis, the Rocket Core was undoubtedly the correct choice. While \acrshort{boom} offers superior standalone \acrshort{cpu} performance, our project's primary goal was to design an \acrshort{soc}, integrate, and evaluate a specialized \acrshort{dnn} accelerator (Gemmini). The vast majority of the performance uplift for our target workload was expected to come from the accelerator, not the host \acrshort{cpu}. Therefore, the selection criteria for the \acrshort{cpu} core prioritized stability, reliability, and ease of integration over raw \acrshort{cpu} performance.

\section{Research Objectives and Contributions}

This thesis aims to demonstrate the capabilities of the Chipyard framework through the implementation of a practical \acrshort{riscv} \acrshort{soc} that integrates a Rocket processor core with the Gemmini \acrshort{dnn} accelerator, while establishing a comprehensive foundation for \acrshort{ceslab}'s transition toward chip prototyping and \acrshort{vlsi}-based development. The specific objectives of this work include:

\begin{enumerate}
    \item \textbf{Framework Evaluation and Migration}: Comprehensively evaluate the Chipyard framework as a replacement for \acrshort{intel} \acrshort{quartus}-based design flows, specifically assessing its \acrshort{vlsi} capabilities, modern interconnect support \allowbreak(Tile\allowbreak Link/\acrshort{axi4}), and suitability for \acrshort{asic} development workflows.
    
    \item \textbf{System Architecture Design}: Develop a comprehensive understanding of the Chipyard framework and utilize it to configure a custom \acrshort{soc} featuring a 64-bit Rocket core coupled with a 16×16 Gemmini systolic array accelerator (Rocket64b1gem16jtag configuration), demonstrating the framework's flexibility and configurability.
    
    \item \textbf{Full-Stack Implementation}: Implement the complete system stack, including hardware generation, software compilation, bootloader configuration, and \acrshort{linux} operating system integration, demonstrating the end-to-end capability of the open-source toolchain from high-level design to functional system.
    
    \item \textbf{\acrshort{fpga} Prototyping and Validation}: Successfully deploy and validate the designed \acrshort{soc} on a Xilinx \acrshort{vc707} \acrshort{fpga} board, providing a realistic hardware platform for system evaluation and serving as a crucial stepping stone toward future \acrshort{asic} implementation.
    
    \item \textbf{Accelerator Integration and Testing}: Demonstrate the integration of the Gemmini accelerator as a \acrfull{rocc}, validate its functionality through comprehensive testing, and showcase how specialized computational units can be seamlessly incorporated into \acrshort{riscv} systems using modern interconnect standards.
    
    \item \textbf{Documentation and Knowledge Transfer}: Create comprehensive documentation of the design process, implementation challenges, and solutions to serve as a foundational reference for future \acrshort{ceslab} projects targeting chip prototyping and \acrshort{vlsi} development.
\end{enumerate}

\section{Significance and Impact}

This work represents the first comprehensive exploration of the \acrshort{riscv} ecosystem within \acrshort{ceslab} and, more importantly, establishes a critical foundation for our laboratory's strategic transition toward chip prototyping and semiconductor fabrication. By successfully implementing a complete, functioning system that combines a \acrshort{riscv} processor with a state-of-the-art \acrshort{dnn} accelerator using industry-standard design flows, this thesis demonstrates the maturity and practicality of open-source hardware design methodologies for serious \acrshort{vlsi} development.

The strategic significance of this work extends far beyond the immediate technical achievements. This thesis serves as a comprehensive feasibility study and implementation roadmap for \acrshort{ceslab}'s evolution from \acrshort{fpga}-based prototyping to actual chip development. The successful demonstration of the Chipyard framework's \acrshort{vlsi} capabilities, coupled with its support for modern interconnect standards and integration with industry-standard \acrshort{eda} tools, validates our laboratory's decision to adopt this platform as the foundation for future chip development projects.

Furthermore, this work establishes a new research paradigm within our laboratory that emphasizes open-source collaboration, reproducible research, and the development of freely available intellectual property. This approach not only reduces the barriers to advanced hardware research but also positions \acrshort{ceslab} as a contributor to the global open-source hardware movement, potentially leading to collaborative opportunities with industry and academic partners pursuing similar chip development goals.

The comprehensive documentation generated by this thesis—covering everything from framework setup and configuration to \acrshort{fpga} deployment and system validation—serves as an invaluable knowledge base for future \acrshort{ceslab} researchers and students. This documentation addresses the practical challenges of implementing complex \acrshort{soc} designs using open-source tools, providing detailed solutions to common problems and establishing best practices for our laboratory's future chip development projects.

Most critically, the full-stack nature of this implementation provides valuable insights into the complex interactions between hardware accelerators, operating systems, and application software—interactions that are often overlooked in studies that focus solely on isolated hardware components. These insights are crucial for developing practical, deployable \acrshort{ai} systems that can deliver their theoretical performance benefits in real-world environments, and they will inform our laboratory's approach to designing optimized \acrshort{asic} implementations in future projects. 

\section{Thesis Structure and Organization}

This thesis is organized into nine chapters that systematically present the foundational argument, technical background, design, implementation, and evaluation of the \acrshort{riscv} \acrshort{soc} with a Gemmini \acrshort{dnn} accelerator.

\textbf{\chapref{chap:riscv_microprocessor_isa}} provides the necessary technical background for this thesis. It begins with a comprehensive overview of the \acrshort{riscv} \acrshort{isa} and its key architectural features. It then presents a detailed description of the Chipyard \acrshort{soc} design framework, the central development platform used in this work.

\textbf{\chapref{chap:SingleCoreGemminiRiscVSystem}} introduces the specific architecture of the system designed and implemented in this project. It details the high-level system structure, the role of \acrshort{chisel}/\acrshort{scala} in hardware design, and the overall methodology for constructing our custom \acrshort{soc} within the Chipyard framework.

\textbf{\chapref{chap:dnn_foundations}} explores the fundamental challenges of \acrlong{dnn} acceleration, including the memory wall problem and the principles of systolic array architectures. This chapter provides the theoretical context for understanding why specialized accelerators like Gemmini are necessary for efficient \acrshort{dnn} computation.

\textbf{\chapref{chap:gemmini_overview}} presents a detailed analysis of the Gemmini accelerator's high-level architecture, including its systolic compute engine, memory hierarchy, and integration model within the \acrshort{soc} framework.

\textbf{\chapref{chap:gemmini_internals}} performs a deeper dive into the microarchitectural details of Gemmini. It examines the internal command and control flow, the instruction processing pipeline that enables efficient \acrshort{dnn} computation.

\textbf{\chapref{chap:implementation}} documents the complete implementation process, from \acrshort{soc} configuration, integration of all system components, and hardware generation to \acrshort{fpga} deployment and system validation. This chapter provides practical insights into the entire development workflow and demonstrates the full-stack, open-source design process.

\textbf{\chapref{chap:results}} presents the results of the project. It includes an analysis of hardware resource utilization, validates the successful boot of a full \acrshort{linux} operating system, and provides a quantitative performance evaluation of the Gemmini accelerator compared to a software-only implementation.

\textbf{\chapref{chap:conclusion_and_future_work}} summarizes the key achievements of this thesis, evaluates the success of the implementation against the stated objectives, acknowledges the limitations of this work, and outlines potential directions for future research and development.

This structured approach ensures a comprehensive understanding of both the theoretical principles underlying modern \acrshort{soc} design and the practical challenges involved in implementing complex, multi-component systems using open-source tools and methodologies.

